\section{Matrices y sus propiedades}
\subsection{Tipos de matrices}

Veamos algunos tipos de matrices que nos encontraremos a lo largo del curso:
\begin{definición}[Matrices introductorias]
    \begin{itemize}
        \item \textbf{Matriz fila:} Una matriz que tiene una sola fila. Por ejemplo, la matriz \( A = [a_1, a_2, a_3] \) es una matriz fila de tamaño \( 1 \times 3 \).
        \item \textbf{Matriz columna:} Una matriz que tiene una sola columna. Por ejemplo, la matriz \( B = \begin{bmatrix} b_1 \\ b_2 \\ b_3 \end{bmatrix} \) es una matriz columna de tamaño \( 3 \times 1 \).
        \item \textbf{Matriz cuadrada:} Una matriz que tiene el mismo número de filas y columnas. Por ejemplo, la matriz \( C = \begin{bmatrix} c_{11} & c_{12} \\ c_{21} & c_{22} \end{bmatrix} \) es una matriz cuadrada de tamaño \( 2 \times 2 \).
        \item \textbf{Matriz diagonal:} Una matriz cuadrada en la que todos los elementos fuera de la diagonal principal son cero. Por ejemplo, la matriz \( D = \begin{bmatrix} d_{11} & 0 \\ 0 & d_{22} \end{bmatrix} \) es una matriz diagonal.
        \item \textbf{Matriz identidad:} Una matriz diagonal en la que todos los elementos de la diagonal principal son uno. Por ejemplo, la matriz \( I = \begin{bmatrix} 1 & 0 \\ 0 & 1 \end{bmatrix} \) es una matriz identidad de tamaño \( 2 \times 2 \).
        \item \textbf{Matriz nula:} Una matriz en la que todos los elementos son cero. Por ejemplo, la matriz \( O = \begin{bmatrix} 0 & 0 \\ 0 & 0 \end{bmatrix} \) es una matriz nula de tamaño \( 2 \times 2 \).
        \item \textbf{Matriz transpuesta:} La matriz obtenida al intercambiar las filas por las columnas de una matriz dada. Si \( A \) es una matriz de tamaño \( m \times n \), su transpuesta \( A^T \) será de tamaño \( n \times m \). Por ejemplo, si \( A = \begin{bmatrix} a_{11} & a_{12} \\ a_{21} & a_{22} \end{bmatrix} \), entonces su transpuesta es \( A^T = \begin{bmatrix} a_{11} & a_{21} \\ a_{12} & a_{22} \end{bmatrix} \).
        \item \textbf{Matriz simétrica:} Una matriz cuadrada que es igual a su transpuesta. Es decir, \( A = A^T \). Por ejemplo, la matriz \( S = \begin{bmatrix} s_{11} & s_{12} \\ s_{12} & s_{22} \end{bmatrix} \) es simétrica.
        \item \textbf{Matriz antisimétrica:} Una matriz cuadrada que es igual a la negativa de su transpuesta. Es decir, \( A = -A^T \). Por ejemplo, la matriz \( T = \begin{bmatrix} 0 & t_{12} \\ -t_{12} & 0 \end{bmatrix} \) es antisimétrica.
        \item \textbf{Matriz triangular:} Una matriz cuadrada que tiene todos los elementos por encima (o por debajo) de la diagonal principal iguales a cero. Si los elementos por debajo de la diagonal son cero, se llama matriz triangular superior; si los elementos por encima son cero, se llama matriz triangular inferior. Por ejemplo, la matriz \( U = \begin{bmatrix} u_{11} & u_{12} \\ 0 & u_{22} \end{bmatrix} \) es triangular superior.
        \item \textbf{Matriz ortogonal:} Una matriz cuadrada \( Q \) es ortogonal si su transpuesta es igual a su inversa, es decir, \( Q^T = Q^{-1} \). Esto implica que \( Q Q^T = Q^T Q = I \), donde \( I \) es la matriz identidad. Tambien se caracteriza porque sus columnas (y filas) forman un conjunto ortonormal en el espacio euclídeo. Por ejemplo, la matriz \( Q = \begin{bmatrix} \frac{1}{\sqrt{2}} & -\frac{1}{\sqrt{2}} \\ \frac{1}{\sqrt{2}} & \frac{1}{\sqrt{2}} \end{bmatrix} \) es ortogonal.
        \item \textbf{Traspuesto conjugado:} Denotado como \( A^* \), es la matriz obtenida al tomar la transpuesta de una matriz y luego tomar el conjugado complejo de cada uno de sus elementos. Por ejemplo, si \( A = \begin{bmatrix} a_{11} & a_{12} \\ a_{21} & a_{22} \end{bmatrix} \), entonces su traspuesto conjugado es \( A^* = \begin{bmatrix} \overline{a_{11}} & \overline{a_{21}} \\ \overline{a_{12}} & \overline{a_{22}} \end{bmatrix} \), donde \( \overline{a_{ij}} \) denota el conjugado complejo de \( a_{ij} \).
        \item \textbf{Matriz definida positiva/negativa:} Una matriz cuadrada \( A \) es definida positiva si para cualquier vector no nulo \( x \), se cumple que \( x^T A x > 0 \). De manera similar, \( A \) es definida negativa si para cualquier vector no nulo \( x \), se cumple que \( x^T A x < 0 \).
        \item \textbf{Matriz semidefinida positiva/negativa:} Una matriz cuadrada \( A \) es semidefinida positiva si para cualquier vector \( x \), se cumple que \( x^T A x \geq 0 \). De manera similar, \( A \) es semidefinida negativa si para cualquier vector \( x \), se cumple que \( x^T A x \leq 0 \).
        \item \textbf{Matriz indefinida:} Una matriz cuadrada \( A \) es indefinida si existen vectores \( x \) y \( y \) tales que \( x^T A x > 0 \) y \( y^T A y < 0 \).
    \end{itemize}
\end{definición}

Ahora definamos algunos tipos de matrices más avanzados:
\begin{definición}[Matriz hermitiana]
    Una matriz cuadrada \( A \) es hermitiana si es igual a su traspuesto conjugado, es decir, \( A = A^* \).\\
    En el caso de matrices con entradas reales, una matriz hermitiana es simplemente una matriz simétrica, es decir:
    \[
        A \in \mathbb{R}^{n \times n} \text{ es hermitiana } \iff A = A^T
    \]
\end{definición}

Veamos un par de ejemplos de matrices hermitianas:
\ejemplo{
    \begin{itemize}
        \item La matriz \( A = \begin{bmatrix} 2 & 3 + i \\ 3 - i & 4 \end{bmatrix} \) es hermitiana, ya que su traspuesto conjugado es \( A^* = \begin{bmatrix} 2 & 3 - i \\ 3 + i & 4 \end{bmatrix} \), que es igual a \( A \).
        \item La matriz \( B = \begin{bmatrix} 1 & 2 \\ 2 & 1 \end{bmatrix} \) es hermitiana, ya que su traspuesto conjugado es \( B^* = \begin{bmatrix} 1 & 2 \\ 2 & 1 \end{bmatrix} \), que es igual a \( B \).
    \end{itemize}
}
    
\begin{definición}[Matriz unitaria]
    Una matriz cuadrada \( U \) es unitaria si su traspuesto conjugado es igual a su inversa, es decir, \( U^* = U^{-1} \). Esto implica que \( U U^* = U^* U = I \), donde \( I \) es la matriz identidad.
\end{definición}

Veamos un par de ejemplos de matrices unitarias:
\ejemplo{
    \begin{itemize}
        \item La matriz \( U = \begin{bmatrix} \frac{1}{\sqrt{2}} & -\frac{1}{\sqrt{2}} \\ \frac{1}{\sqrt{2}} & \frac{1}{\sqrt{2}} \end{bmatrix} \) es unitaria, ya que su traspuesto conjugado es \( U^* = \begin{bmatrix} \frac{1}{\sqrt{2}} & \frac{1}{\sqrt{2}} \\ -\frac{1}{\sqrt{2}} & \frac{1}{\sqrt{2}} \end{bmatrix} \), y se cumple que \( U U^* = I \).
        \item La matriz \( V = \begin{bmatrix} 0 & 1 \\ 1 & 0 \end{bmatrix} \) es unitaria, ya que su traspuesto conjugado es \( V^* = \begin{bmatrix} 0 & 1 \\ 1 & 0 \end{bmatrix} \), y se cumple que \( V V^* = I \).
    \end{itemize}
}

\begin{observación}
    Se puede apreciar que si $U$ es una matriz unitaria compleja:
    \[
        U^* = U^{-1} \implies 
        U^* U = I \implies
        U^* U =
        \begin{pmatrix}
            u_1^* \\
            \vdots \\
            u_n^* \\
        \end{pmatrix}
        \begin{pmatrix}
            u_1 & \dots & u_n \\
        \end{pmatrix} =
        \begin{pmatrix}
            u_1^* u_1 & \dots & u_1^* u_n \\
            \vdots & \ddots & \vdots \\
            u_n^* u_1 & \dots & u_n^* u_n
        \end{pmatrix} =
        \begin{pmatrix}
            1 & \dots & 0 \\
            \vdots & \ddots & \vdots \\
            0 & \dots & 1
        \end{pmatrix}
    \]
    De lo cual se deduce que los vectores columna de $U$ forman un conjunto ortonormal en el espacio complejo, es decir, $u_i^* u_j = \delta_{ij}$, ya que el producto $u_i^* u_j$ representa el producto interno entre los vectores $u_i$ y $u_j$ en el espacio complejo, y hemos visto que $u_i^* u_j = 0$ si $i \neq j$ y $u_i^* u_i = ||u_i||^2 = 1$.
\end{observación}

\begin{definición}[Matriz normal]
    Una matriz cuadrada \( A \) es normal si conmuta con su traspuesto conjugado, es decir, si se cumple que \( A A^* = A^* A \).
\end{definición}

Veamos un par de ejemplos de matrices normales:
\ejemplo{
    \begin{itemize}
        \item La matriz \( A = \begin{bmatrix} 2 & 0 \\ 0 & 3 \end{bmatrix} \) es normal, ya que su traspuesto conjugado es \( A^* = \begin{bmatrix} 2 & 0 \\ 0 & 3 \end{bmatrix} \), y se cumple que \( A A^* = A^* A = \begin{bmatrix} 4 & 0 \\ 0 & 9 \end{bmatrix} \).
        \item La matriz \( B = \begin{bmatrix} 1 & i \\ -i & 1 \end{bmatrix} \) es normal, ya que su traspuesto conjugado es \( B^* = \begin{bmatrix} 1 & i \\ -i & 1 \end{bmatrix} \), y se cumple que \( B B^* = B^* B = \begin{bmatrix} 2 & 0 \\ 0 & 2 \end{bmatrix} \).
    \end{itemize}
}

\begin{observación}
    Notese que si $A$ es hermitiana, entonces $A = A^*$, por lo que:
    \[
        A A^* = A A = A^2 = A^* A
    \]
    Por lo tanto, toda matriz hermitiana es normal.\\
    Algo similar ocurre con las matrices unitarias, ya que si $U$ es unitaria, entonces $U^* = U^{-1}$, por lo que:
    \[
        U U^* = U U^{-1} = I = U^{-1} U = U^* U
    \]
    Por lo tanto, toda matriz unitaria es normal.
\end{observación}

\begin{proposición}[Porducto de matrices triangulares]
    Sean $A, B \in \mathbb{C}^{n \times n}$ dos matrices triangulares inferiores (superiores). Entonces, el producto \( C = AB \) es también una matriz triangular inferior (superior).
\end{proposición}

\begin{proof}
    Podemos aplicar una induccion sencilla sobre el tamaño de las matrices.\\
    \textbf{Caso base:} Para \( n = 1 \), las matrices son simplemente escalares, y el producto de dos escalares es un escalar, que puede considerarse como una matriz triangular inferior (superior) trivialmente.\\
    \textbf{Paso inductivo:} Supongamos que la proposición es cierta para matrices de tamaño \( (n-1) \times (n-1) \). Consideremos dos matrices triangulares inferiores \( A, B \in \mathbb{C}^{n \times n} \):
    \[
    A =
    \begin{array}{c|c}
        a_11 & 0 \\
        \hline
        \alpha & A_{(n-1)}
    \end{array}
    , \quad
    B =
    \begin{array}{c|c}
        b_11 & 0 \\
        \hline
        \beta & B_{(n-1)}
    \end{array}
    \]
    Y su producto:
    \[
    C = AB =
    \begin{array}{c|c}
        a_{11} b_{11} & 0 \\
        \hline
        \alpha b_{11} + A_{(n-1)} \beta & A_{(n-1)} B_{(n-1)}
    \end{array}
    \]
    Por la hipótesis inductiva, \( A_{(n-1)} B_{(n-1)} \) es una matriz triangular inferior de tamaño \( (n-1) \times (n-1) \). Además, el elemento en la posición \( (n, 1) \) de \( C \) es \( \alpha b_{11} + A_{(n-1)} \beta \), que es un vector columna. Por lo tanto, todos los elementos por encima de la diagonal principal de \( C \) son cero, lo que demuestra que \( C \) es una matriz triangular inferior.\\
    Por lo tanto, por el principio de inducción matemática, la proposición es cierta para todas las matrices triangulares inferiores (superiores) de tamaño \( n \times n \).
\end{proof}

\begin{definición}[Matriz banda]
    Se define una matriz banda como una matriz cuadrada en la que los elementos no nulos se encuentran dentro de una franja diagonal alrededor de la diagonal principal. Esta franja está determinada por dos enteros no negativos \( k_1 \) y \( k_2 \), que representan el número de subdiagonales y superdiagonales, respectivamente. Formalmente, una matriz \( A \in \mathbb{R}^{n \times n} \) es una matriz banda si:
    \[
        a_{ij} = 0 \quad \text{si } |i - j| > k
    \]
    donde \( k = \max(k_1, k_2) \).\\
\end{definición}

Veamos un par de ejemplos de matrices banda:
\ejemplo{
    \begin{itemize}
        \item Una matriz tridiagonal es una matriz banda con \( k_1 = k_2 = 1 \). Por ejemplo, la matriz \( A = \begin{bmatrix} a_{11} & a_{12} & 0 & 0 \\ a_{21} & a_{22} & a_{23} & 0 \\ 0 & a_{32} & a_{33} & a_{34} \\ 0 & 0 & a_{43} & a_{44} \end{bmatrix} \) es tridiagonal.
        \item Una matriz pentadiagonal es una matriz banda con \( k_1 = k_2 = 2 \). Por ejemplo, la matriz \( B = \begin{bmatrix} b_{11} & b_{12} & b_{13} & 0 & 0 \\ b_{21} & b_{22} & b_{23} & b_{24} & 0 \\ b_{31} & b_{32} & b_{33} & b_{34} & b_{35} \\ 0 & b_{42} & b_{43} & b_{44} & b_{45} \\ 0 & 0 & b_{53} & b_{54} & b_{55} \end{bmatrix} \) es pentadiagonal.
    \end{itemize}
}

\begin{definición}[Matriz de diagonal estrictamente dominante]
    Una matriz cuadrada \( A = [a_{ij}] \in \mathbb{R}^{n \times n} \) se dice que es de diagonal estrictamente dominante si para cada fila \( i \) se cumple que:
    \[
        |a_{ii}| > \sum_{\substack{j=1 \\ j \neq i}}^{n} |a_{ij}|
    \]
    Es decir, el valor absoluto del elemento en la diagonal principal de cada fila es mayor que la suma de los valores absolutos de los demás elementos en esa fila.
\end{definición}

Veamos un par de ejemplos de matrices de diagonal estrictamente dominante:
\ejemplo{
    \begin{itemize}
        \item La matriz \( A = \begin{bmatrix} 4 & 1 & 2 \\ 3 & 5 & 1 \\ 1 & 1 & 6 \end{bmatrix} \) es de diagonal estrictamente dominante, ya que:
        \[
            |4| > |1| + |2|, \quad |5| > |3| + |1|, \quad |6| > |1| + |1|
        \]
        \item La matriz \( B = \begin{bmatrix} 10 & 2 & 1 \\ 1 & 8 & 2 \\ 2 & 1 & 9 \end{bmatrix} \) es de diagonal estrictamente dominante, ya que:
        \[
            |10| > |2| + |1|, \quad |8| > |1| + |2|, \quad |9| > |2| + |1|
        \]
    \end{itemize}
}

\begin{definición}[Traza de una matriz]
    La traza de una matriz cuadrada \( A = [a_{ij}] \in \mathbb{R}^{n \times n} \) se define como la suma de los elementos en la diagonal principal de la matriz. Matemáticamente, se expresa como:
    \[
        \text{tr}(A) = \sum_{i=1}^{n} a_{ii}
    \]
    donde \( a_{ii} \) son los elementos en la diagonal principal de la matriz \( A \).    
\end{definición}

Veamos un par de ejemplos de cálculo de la traza de una matriz:
\ejemplo{
    \begin{itemize}
        \item Para la matriz \( A = \begin{bmatrix} 1 & 2 & 3 \\ 4 & 5 & 6 \\ 7 & 8 & 9 \end{bmatrix} \), la traza es:
        \[
            \text{tr}(A) = 1 + 5 + 9 = 15
        \]
        \item Para la matriz \( B = \begin{bmatrix} 2 & 0 & 1 \\ 0 & 3 & 4 \\ 5 & 6 & 7 \end{bmatrix} \), la traza es:
        \[
            \text{tr}(B) = 2 + 3 + 7 = 12
        \]
    \end{itemize}
}

\begin{proposición}[Propiedades de la traza]
    Veamos algunas propiedades importantes de la traza de una matriz:
    \begin{itemize}
        \item \textbf{Linealidad:} La traza es una función lineal, es decir, para dos matrices cuadradas \( A \) y \( B \) del mismo tamaño y un escalar \( c \), se cumple que:
        \[
            \text{tr}(A + B) = \text{tr}(A) + \text{tr}(B)
        \]
        \[
            \text{tr}(cA) = c \cdot \text{tr}(A)
        \]
        \item \textbf{Invarianza bajo trasposición:} La traza de una matriz es igual a la traza de su transpuesta:
        \[
            \text{tr}(A) = \text{tr}(A^T)
        \]
        \item \textbf{Invarianza bajo conjugación:} La traza de una matriz es igual a la traza de su traspuesto conjugado:
        \[
            \text{tr}(\overline{A}) = \text{tr}(A^*)
        \]
        \item \textbf{Cyclicidad:} La traza del producto de dos matrices es invariante bajo el cambio de orden:
        \[
            \text{tr}(AB) = \text{tr}(BA)
        \]
    \end{itemize}
    Las propiedades anteriores se prueban directamente a partir de la definición de traza y las propiedades básicas de las operaciones matriciales, veamos solo la ultima:
    \[
        C = AB \implies
        c_{ii} = \sum_{k=1}^{n} a_{ik} b_{ki} \implies
        \text{tr}(C) = \sum_{i=1}^{n} c_{ii} = 
        \sum_{i=1}^{n} \sum_{k=1}^{n} a_{ik} b_{ki}
    \]
    \[
        D = BA \implies
        d_{ii} = \sum_{k=1}^{n} b_{ik} a_{ki} \implies
        \text{tr}(D) = \sum_{i=1}^{n} d_{ii} = 
        \sum_{i=1}^{n} \sum_{k=1}^{n} b_{ik} a_{ki}
    \]
    Cambiando el orden de las sumas en ambos casos, se llega a la igualdad:
    \[
        \text{tr}(C) = \text{tr}(D)
    \]
\end{proposición}

\begin{definición}[Rango de una matriz]
    El rango de una matriz \( A \in \mathbb{R}^{m \times n} \) es el número máximo de columnas linealmente independientes que tiene la matriz. Se denota como \( \text{rg}(A) \). El rango también puede ser interpretado como la dimensión del espacio generado por las columnas de la matriz, o como el número de filas no nulas en la forma escalonada reducida por filas de la matriz.
\end{definición}

\begin{observación}
    Notese que el rango de una matriz \( A \) es igual al rango de su transpuesta \( A^T \), es decir, \( \text{rg}(A) = \text{rg}(A^T) \). Esto se debe a que el número de columnas linealmente independientes en \( A \) es igual al número de filas linealmente independientes en \( A^T \), ya que las filas de \( A^T \) corresponden a las columnas de \( A \).
\end{observación}

\begin{definición}[Determinante]
    El determinante de una matriz cuadrada \( A \in \mathbb{R}^{n \times n} \) es un valor escalar que se asocia a la matriz y que proporciona información sobre ciertas propiedades de la misma, como su invertibilidad y el volumen de transformación lineal que representa. El determinante se denota como \( \det(A) \) o \( |A| \), y se calcula como:
    \[
        \det(A) = \sum_{\sigma \in S_n} \text{sgn}(\sigma) \prod_{i=1}^{n} a_{i, \sigma(i)}
    \]
    donde \( S_n \) es el conjunto de todas las permutaciones de los números \( 1, 2, \ldots, n \), y \( \text{sgn}(\sigma) \) es el signo de la permutación \( \sigma \).
\end{definición}

\begin{proposición}[Propiedades del determinante]
    Veamos algunas propiedades importantes del determinante de una matriz:
    \begin{itemize}
        \item \textbf{Multiplicatividad:} El determinante del producto de dos matrices es igual al producto de sus determinantes:
        \[
            \det(AB) = \det(A) \cdot \det(B)
        \]
        \item \textbf{Determinante de matrices diagonales:} El determinante de una matriz diagonal es igual al producto de los elementos en la diagonal principal:
        \[
            \det(D) = \prod_{i=1}^{n} d_{ii}
        \]
        \item \textbf{Transposición:} El determinante de una matriz es igual al determinante de su transpuesta:
        \[
            \det(A) = \det(A^T)
        \]
        \item \textbf{Escalar:} Si \( c \) es un escalar y \( A \) es una matriz cuadrada de tamaño \( n \times n \), entonces:
        \[
            \det(cA) = c^n \cdot \det(A)
        \]
        \item \textbf{Inversión:} Si \( A \) es una matriz invertible, entonces:
        \[
            \det(A^{-1}) = \frac{1}{\det(A)}
        \]
    \end{itemize}
    Todas estas propiedades se pueden demostrar utilizando la definición del determinante y las propiedades básicas de las operaciones matriciales.
\end{proposición}

\begin{proposición}[Determinante de una matriz unitaria]
    Sea \( U \in \mathbb{C}^{n \times n} \) una matriz unitaria. Entonces, el valor absoluto de su determinante es igual a 1, es decir:
    \[
        |\det(U)| = 1
    \]
\end{proposición}

\begin{proof}
    Sea \( U \) una matriz unitaria, por definición, se cumple que \( U^* = U^{-1} \). Aplicando la propiedad de multiplicatividad del determinante, tenemos:
    \[
    U U^* = I \implies
    \det(U U^*) = \det(I) = 1 \implies 
    \det(U) \cdot \det(U^*) =
    \det(U) \cdot \overline{\det(U)} =
    |\det(U)|^2 = 1
    \]
    De donde se deduce que:
    \[
        |\det(U)| = 1
    \]
\end{proof}

\begin{teorema}[Caracterización de las matrices invertibles]
    Sea \( A \in \mathbb{R}^{n \times n} \) una matriz cuadrada. Entonces, las siguientes afirmaciones son equivalentes:
    \begin{enumerate}
        \item \( A \) es invertible, es decir $\exists A^{-1}$ tal que \( AA^{-1} = A^{-1}A = I \).
        \item El determinante de \( A \) es diferente de cero, es decir, \( \det(A) \neq 0 \).
        \item El espectro de \( A \) no contiene el valor cero, es decir, \( 0 \notin \sigma(A) \).
    \end{enumerate}
\end{teorema}

\begin{proof}
    Veamos las implicaciones:
    \begin{itemize}
        \item[$2) \iff 3)$] Sabemos que el espectro de una matriz \( A \) está dado por los valores propios de la matriz, que son las soluciones de la ecuación característica \( \det(A - \lambda I) = 0 \). Si \( 0 \in \sigma(A) \), entonces existe un valor propio \( \lambda = 0 \) tal que \( \det(A - 0 \cdot I) = \det(A) = 0 \). Por lo tanto, si \( \det(A) \neq 0 \), entonces \( 0 \notin \sigma(A) \), y viceversa.
        \item[$1) \implies 2)$] Si \( A \) es invertible, entonces existe una matriz \( A^{-1} \) tal que \( AA^{-1} = I \). Aplicando la propiedad de multiplicatividad del determinante, tenemos:
        \[
            \det(AA^{-1}) = \det(I) = 1 \implies
            \det(A) \cdot \det(A^{-1}) = 1 \implies
            \det(A) = (\det(A^{-1}))^{-1} \neq 0
        \]
        De donde se deduce que \( \det(A) \neq 0 \).
        \item[$2) \implies 1)$] Si \( \det(A) \neq 0 \), entonces la ecuación característica \( \det(A - \lambda I) = 0 \) no tiene \( \lambda = 0 \) como solución, lo que implica que \( A \) no tiene un núcleo no trivial. Por lo tanto, la transformación lineal asociada a \( A \) es inyectiva y, dado que \( A \) es cuadrada, también es sobreyectiva. Por el teorema de la inversa, esto implica que \( A \) es invertible.\\
        Otra forma de demostrarlo es utilizando la descomposición LU o la eliminación de Gauss. Si \( \det(A) \neq 0 \), entonces \( A \) puede ser reducida a la matriz identidad mediante operaciones elementales, es decir, $\exists E_1, E_2, \dots, E_k$ matrices elementales tales que:
        \[
            A E =
            A E_1 E_2 \dots E_k = I
        \]
        Donde \( E = E_1 E_2 \dots E_k \) es una matriz invertible, ya que es el producto de matrices elementales. Por lo tanto, podemos escribir:
        \[
            A = E^{-1} I = E^{-1}
        \]
        Y se deduce que \( A \) es invertible, con \( A^{-1} = E \).
    \end{itemize}
\end{proof}

Veamos ahora varias propiedades interesantes del calculo de determinantes.

\begin{corolario}[Determinante de una matriz triangular]
    El determinante de una matriz triangular (superior o inferior) es igual al producto de los elementos en la diagonal principal. Es decir, si$( A $ es una matriz triangular, entonces:
    \[
        \det(A) = \prod_{i=1}^{n} a_{ii}
    \]
    Esto se da ya que si en una permutacion tenemos un elemento fuera de la diagonal principal, entonces el producto de los elementos correspondientes a esa permutacion sera cero, ya que si el elemento esta por debajo de la diagonal principal (en una triangular superior) o por encima de la diagonal principal (en una triangular inferior), entonces ese elemento es cero, lo que hace que el producto de los elementos correspondientes a esa permutacion sea cero, esta en el otro lado de la diagonal, debe haber un elemento fuera de la diagonal principal, lo que hace que el producto de los elementos correspondientes a esa permutacion sea cero, y asi sucesivamente para todas las permutaciones que tengan un elemento fuera de la diagonal principal, lo que hace que el unico producto no nulo sea el producto de los elementos en la diagonal principal, que corresponde a la permutacion identidad.
\end{corolario}

\begin{corolario}
    Sea $M = \left( \begin{array}{c|c} A & 0 \\ \hline 0 & B \end{array} \right)$ una matriz con "solo 2 esquinas", entonces ocurre lo siguiente:
    \[
        \det(M) =
        \begin{array}{|c|c|}
            A & 0 \\
            \hline
            0 & B
        \end{array} =
        \begin{array}{|c|c|}
            A & 0 \\
            \hline
            0 & I
        \end{array}
        \cdot
        \begin{array}{|c|c|}
            I & 0 \\
            \hline
            0 & B
        \end{array} =
        \det(A) \cdot \det(B)
    \]
\end{corolario}

\begin{corolario}
    Sea $M = \left( \begin{array}{c|c} A & C \\ \hline 0 & B \end{array} \right)$ una matriz con "solo 3 esquinas", entonces ocurre lo siguiente:
    \[
        \begin{array}{|c|c|}
            I & \Gamma \\
            \hline
            0 & I
        \end{array}
        \cdot
        \begin{array}{|c|c|}
            A & C \\
            \hline
            0 & B
        \end{array} =
        \begin{array}{|c|c|}
            A & C + \Gamma B \\
            \hline
            0 & B
        \end{array}
    \]
    Ahora si tomamos $\Gamma = -CB^{-1}$, entonces:
    \[
        \begin{array}{|c|c|}
            I & -CB^{-1} \\
            \hline
            0 & I
        \end{array}
        \cdot
        \begin{array}{|c|c|}
            A & C \\
            \hline
            0 & B
        \end{array} =
        \begin{array}{|c|c|}
            A & 0 \\
            \hline
            0 & B
        \end{array}
    \]
    Esto lo podemos hacer ya que si $B$ no es invertible, entonces habria una combinación lineal de las filas de $B$ que daria como resultado la fila nula, la cual al replicarse en la esquina inferior derecha de $M$ haria que $M$ no sea invertible, lo que implicaria que $\det(M) = 0$, y por lo tanto, $\det(A) \cdot \det(B) = 0$, lo que implica que $\det(B) = 0$.\\
    Por lo tanto, si $B$ es invertible, entonces podemos realizar la operación anterior para obtener una matriz con "solo 2 esquinas", y aplicar el resultado del corolario anterior para concluir que:
    \[
        \det(M) = \det(A) \cdot \det(B)
    \]
    Ya que el primer determinante $\begin{array}{|c|c|}
            I & \Gamma \\
            \hline
            0 & I
        \end{array} $ es igual a 1, por ser una matriz triangular superior con 1's en la diagonal principal.
\end{corolario}

\subsection{Autovalores y autovectores}

Para definir los autovalores y autovectores de una matriz, primero necesitamos definir el concepto de espectro de una matriz, y el polinomio característico de una matriz.\\

\begin{definición}[Espectro de una matriz]
    El espectro de una matriz cuadrada $ A \in \mathbb{R}^{n \times n} $ es el conjunto de todos los valores propios de la matriz, es decir, el conjunto de todos los escalares $ \lambda $ para los cuales existe un vector no nulo $ v $ tal que:
    \[
        Av = \lambda v
    \]
    El espectro se denota como $\sigma(A) = sp(A)$.
\end{definición}

\begin{definición}[Polinomio característico de una matriz]
    El polinomio característico de una matriz cuadrada $ A \in \mathbb{R}^{n \times n} $ es un polinomio en la variable $ \lambda $ definido como:
    \[
        p_A(\lambda) = \det(A - \lambda I)
    \]
    donde $ I $ es la matriz identidad de tamaño $ n \times n $.\\
    Los valores propios de la matriz $ A $ son las raíces del polinomio característico, es decir, los valores de $ \lambda $ para los cuales $ p_A(\lambda) = 0 $.
\end{definición}

Ahora sí, podemos definir formalmente los autovalores y autovectores de una matriz:

\begin{definición}[Autovalores y autovectores de una matriz]
    Dada una matriz cuadrada $ A \in \mathbb{R}^{n \times n} $, un escalar $ \lambda $ se llama autovalor de $ A $ si existe un vector no nulo $ v \in \mathbb{R}^n $ tal que:
    \[
        Av = \lambda v
    \]
    En este caso, el vector $ v $ se llama autovector asociado al autovalor $ \lambda $.\\
    El conjunto de todos los autovectores asociados a un autovalor $ \lambda $ forma un subespacio vectorial llamado espacio propio o eigenspace correspondiente a $ \lambda $. Este espacio se denota como $ E_\lambda = \{ v \in \mathbb{R}^n : Av = \lambda v \} $.
\end{definición}

Notese que a los autovalores también se les conoce como valores propios, y a los autovectores como vectores propios.\\

Una forma de calcular los autovalores de una matriz es utilizando el polinomio característico, ya que los autovalores son las raíces de dicho polinomio. Para calcular los autovectores asociados a un autovalor $ \lambda $, se puede resolver el sistema de ecuaciones lineales dado por:
\[
    (A - \lambda I)v = 0
\]
que corresponde a encontrar el núcleo de la matriz $ A - \lambda I $. El espacio propio $ E_\lambda $ es precisamente el núcleo de esta matriz, y sus elementos son los autovectores asociados al autovalor $ \lambda $.\\

Notese también que el espectro es el conjunto de autovalores, y que cada autovalor puede tener uno o más autovectores asociados, dependiendo de su multiplicidad algebraica y geométrica.

\begin{definición}[Multiplicidad algebraica y geométrica de un autovalor]
    Dado un autovalor $ \lambda $ de una matriz $ A \in \mathbb{R}^{n \times n} $, se define la multiplicidad algebraica de $ \lambda $ como el número de veces que $ \lambda $ aparece como raíz del polinomio característico $ p_A(\lambda) $.\\
    Por otro lado, la multiplicidad geométrica de $ \lambda $ se define como la dimensión del espacio propio $ E_\lambda $, es decir, el número máximo de autovectores linealmente independientes asociados al autovalor $ \lambda $.\\
    Es importante destacar que la multiplicidad geométrica de un autovalor siempre es menor o igual a su multiplicidad algebraica.
\end{definición}

\begin{definición}[Radio espectral de una matriz]
    El radio espectral de una matriz cuadrada $ A \in \mathbb{R}^{n \times n} $ se define como el máximo valor absoluto de los autovalores de la matriz. Matemáticamente, se expresa como:
    \[
        \rho(A) = \max_{\lambda \in \sigma(A)} |\lambda|
    \]
    El radio espectral proporciona información sobre el comportamiento asintótico de las potencias de la matriz $ A $, ya que si $ \rho(A) < 1 $, entonces las potencias de $ A $ tienden a cero, mientras que si $ \rho(A) > 1 $, entonces las potencias de $ A $ crecen sin límite.
\end{definición}

Veamos ahora como el espectro se mantiene ante la conmutación del producto de matrices.

\begin{proposición}
    Si $ A, B \in \mathbb{R}^{n \times n} $ son dos matrices cuadradas, entonces el espectro de $ AB $ es igual al espectro de $ BA $, es decir:
    \[
        \sigma(AB) = \sigma(BA)
    \]
\end{proposición}

\begin{proof}
    Para ver esto, veamos como $ \forall \lambda \in \sigma(AB) $, entonces existe un vector no nulo $ v $ tal que:
    \[
        BAv = \lambda v
    \]
    Es decir, que $ \lambda $ es un autovalor de $ BA $ (la implicacion contraria es análoga). Para ello distingamos dos casos:
    \begin{itemize}
        \item[$\lambda = 0$] En este caso:
        \[
            0 \in \sigma(AB) \iff
            \det(AB) = 0 \det(BA) \iff 
            0 \in \sigma(BA)
        \]
        \item[$\lambda \neq 0$] En este caso, como $ v \neq 0 $ es un autovector asociado al autovalor $ \lambda $ de $ AB $, entonces $ ABv = \lambda v $. Multiplicando ambos lados por $ B $, obtenemos:
        \[
            BA(Bv) =
            B(ABv) = 
            \lambda (Bv)
        \]
        Por lo que $ Bv $ es un autovector asociado al autovalor $ \lambda $ de $ BA $. Además, como $ v $ es un vector no nulo y $ \lambda \neq 0 $, entonces $ Bv $ también es un vector no nulo, lo que garantiza que $ \lambda \in \sigma(BA) $, ya que ademas $\det(B) \neq 0$ (si $B$ no es invertible, entonces $0 \in \sigma(B)$, lo que implica que $0 \in \sigma(BA)$, lo que contradice la suposición de que $\lambda \neq 0$).
    \end{itemize}
\end{proof}

Veamos ahora como el espectro de una matriz se traslada al de sus potencias.

\begin{proposición}
    Si $ A \in \mathbb{R}^{n \times n} $ es una matriz cuadrada, entonces el espectro de $ A^k $ es igual al conjunto de los autovalores de $ A $ elevados a la potencia $ k $, es decir:
    \[
        \sigma(A^k) = \{ \lambda^k : \lambda \in \sigma(A) \}
    \]
\end{proposición}

\begin{proof}
    Para ver esto, veamos como $ \forall \lambda \in \sigma(A) $, entonces existe un vector no nulo $ v $ tal que:
    \[
        Av = \lambda v
    \]
    Ahora apliquemos inducción sobre $ k $ para demostrar que $ A^k v = \lambda^k v $:
    \begin{itemize}
        \item \textbf{Caso base:} Para $ k = 1 $, se cumple que $ A^1 v = Av = \lambda v $, lo que es cierto por la definición de autovalor y autovector.
        \item \textbf{Paso inductivo:} Supongamos que para algún $ k \geq 1 $, se cumple que $ A^k v = \lambda^k v $. Entonces, para $ k + 1 $, tenemos:
        \[
            A^{k+1} v =
            A(A^k v) =
            A(\lambda^k v) =
            \lambda^k (Av) =
            \lambda^k (\lambda v) =
            \lambda^{k+1} v
        \]
        Por lo tanto, por el principio de inducción matemática, se cumple que $ A^k v = \lambda^k v $ para todo $ k \geq 1 $.
    \end{itemize}
    De donde se deduce que $ \lambda^k $ es un autovalor de $ A^k $ con el mismo autovector $ v $. Además, como $ v $ es un vector no nulo, entonces $ \lambda^k $ también es un valor propio de $ A^k $. Por lo tanto, se concluye que:
    \[
        \sigma(A^k) = \{ \lambda^k : \lambda \in \sigma(A) \}
    \]
\end{proof}

\subsection{Teoremas importantes}

\begin{teorema}[Descomposición QR]
    Sea $A \in \mathbb{K}^{m \times n}$ una matriz de rango máximo.\\
    Entonces existen una matriz unitaria (ortogonal) $Q \in \mathbb{K}^{m \times m}$ y una matriz triangular superior $R \in \mathbb{K}^{m \times n}$ tal que:
    \[
        A = QR \underbrace{\iff}_{Q^{-1} = Q^*}
        Q^* A = R
    \]
\end{teorema}

Antes de la demostración, conviene definir el operador de proyección ortogonal de un vector sobre otro.

\begin{definición}[Proyección ortogonal de un vector sobre otro]
    Dados dos vectores $ u, a \in \mathbb{K}^n $, la proyección ortogonal de $ a $ sobre $ u $ se define como el vector:
    \[
        \begin{array}{rccl}
            p : & \mathbb{K}^n \times \mathbb{K}^n & \to & \mathbb{K}^n \\
            & (u, a) & \mapsto & p(u, a) = 
            \frac{\langle u, a \rangle}{\|u\|^2} u =
            \langle \frac{u}{\|u\|}, a \rangle \frac{u}{\|u\|}
        \end{array}
    \]
    Teniendo en cuenta que:
    \[
        \langle u, a \rangle = u^* a = a^* u
    \]
    \[
        \|u\|^2 = \langle u, u \rangle = u^* u
    \]
\end{definición}

\begin{proof}
    Vamos a realizar una demostración constructiva utilizando el proceso de Gram-Schmidt para ortogonalizar las columnas de la matriz $ A $, es decir, queremos contruir una base ortogonal tal que $Q = (q_1, q_2, \dots, q_n)$, y para ello haremos lo siguiente:\\
    Iremos tomando vectores $u_i$ que sean ortogonales entre sí, y luego los normalizaremos para obtener los vectores $q_i$ que formarán la matriz unitaria $Q$.\\
    Para el primer vector, tomamos $ u_1 = a_1 $, donde $a_1$ es la primera columna de $A$. Luego, normalizamos $u_1$ para obtener $ q_1 = \frac{u_1}{\|u_1\|} $.\\
    Ahora para obtener el segundo vector, tomamos $ u_2 = a_2 - p(u_1, a_2) $, donde $ p(u_1, a_2) $ es la proyección ortogonal de $ a_2 $ sobre $ u_1 $. Luego, normalizamos $ u_2 $ para obtener $ q_2 = \frac{u_2}{\|u_2\|} $.\\
    Para el tercer vector, tomamos $ u_3 = a_3 - p(u_1, a_3) - p(u_2, a_3) $, donde $ p(u_1, a_3) $ y $ p(u_2, a_3) $ son las proyecciones ortogonales de $ a_3 $ sobre $ u_1 $ y $ u_2 $, respectivamente. Luego, normalizamos $ u_3 $ para obtener $ q_3 = \frac{u_3}{\|u_3\|} $.\\
    Continuando de esta manera, para el $ k $-ésimo vector, tomamos:
    \[
        u_k = a_k - \sum_{j=1}^{k-1} p(u_j, a_k)
    \]
    Luego, normalizamos $ u_k $ para obtener $ q_k = \frac{u_k}{\|u_k\|} $.\\
    De esta forma, obtenemos una base ortonormal $ Q = (q_1, q_2, \dots, q_n) $ que forma una matriz unitaria.\\
    Todo esto se puede entender geometricamente como ir proyectando cada columna de $ A $ sobre el espacio generado por las columnas anteriores, quitandole la parte que esta en ese espacio, y luego normalizando el resultado para obtener un vector unitario, vease el siguiente dibujo:
    \begin{center}
        \begin{tikzpicture}
            % Coordenadas
            % Inicial
            \coordinate (O) at (0, 0);
            % Final de q_1
            \coordinate (Q1) at (3, 1);
            % Final de a
            \coordinate (A) at (1, 2);

            % Dibujo de los vectores
            \draw[->, thick, red] (O) -- (Q1) node[below] {$q_1$};
            \draw[->, thick, blue] (O) -- (A) node[midway, above] {$a$};
            % Proyección de a sobre u
            \coordinate (P) at ($(Q1)!(A)!(O)$);
            \draw[->, thick, dashed] (O) -- (P) node[below] {$p(q_1, a)$};
            % Vector resultante de u_2
            \draw[->, thick, dotted, orange] (P) -- (A) node[midway, above] {$u_2$};
        \end{tikzpicture}
    \end{center}
    Así, hemos construido la matriz unitaria $ Q $ a partir de las columnas de $ A $, ya que los vectores $ q_i $ son ortonormales y forman una base del espacio generado por las columnas de $ A $, es decir, de $\mathbb{K}^n$.\\
    Gracias a esta construcción, podemos definir la matriz triangular superior $ R $ como:
    \[
        R = Q^* A =
        \begin{pmatrix}
            q_1^* \\
            q_2^* \\
            \vdots \\
            q_n^*
        \end{pmatrix}
        \cdot
        \begin{pmatrix}
            a_1 & a_2 & \cdots & a_n
        \end{pmatrix} =
        \begin{pmatrix}
            \color{blue} q_1^* a_1 & \color{blue}q_1^* a_2 & \cdots & \color{blue}q_1^* a_n \\
            \color{red} 0 & \color{blue} q_2^* a_2 & \cdots & \color{blue} q_2^* a_n \\
            \vdots & \vdots & \ddots & \vdots \\
            \color{red} 0 & \color{red} 0 & \cdots & \color{blue} q_n^* a_n
        \end{pmatrix}
    \]
    Estos $0$ aparecen ya que:
    \[
        q_m^* a_k =
        \langle q_m, a_k \rangle =
        \langle q_m, u_k + \sum_{j=1}^{k-1} p(u_j, a_k) \rangle =
        \langle q_m, u_k \rangle + \sum_{j=1}^{k-1} \langle q_m, p(u_j, a_k) \rangle
    \]
    Y ahora:
    \[
        \langle q_m, u_k \rangle =
        \begin{cases}
            0 & \text{si } m < k \\
            \|u_k\| & \text{si } m = k
        \end{cases}
    \]
    Por lo que si 
\end{proof}